% DRAFT 6 — Discussion and Conclusion, split (were merged at 235 words).
% All p-values from the REGENERATED 12-model bootstrap (66 pairs).
% Parameter-efficiency comparisons appear HERE ONCE and nowhere else.

\section{Discussion}
\label{sec:discussion}

\subsection{What a Single Accuracy Number Hides}

The strict and judge-audited regimes rank the same twelve models on the same 406 items, and they
disagree about the ordering. That disagreement is not noise to be averaged away; it is the
measurement telling us that two distinct quantities have been collapsed into one. Strict accuracy
measures what a format-sensitive pipeline can extract. Judge-audited accuracy measures what the model
appears to know. For most models the two are close, which is why the conflation usually goes
unnoticed. For a model that narrates its reasoning, they differ by twenty-three points and eleven ranks.

The practical consequence is that a leaderboard cannot be read without knowing which question it
answers. If the intended use is an automated compliance pipeline that parses model output
mechanically, the strict number is the relevant one and DeepSeek-R1-Distill really is a poor choice. If
the question is which model reasons most reliably about Indian regulatory text, the strict number is
actively misleading about that same model. Benchmarks that report one figure force the reader to
guess which of these they are looking at.

We would not argue for replacing string matching with judge-based scoring. Judges have their own
failure modes, and ours is no exception: the Gemini pilot's 28-item manual audit found every error
running the same direction (Appendix~\ref{app:judge})---over-crediting a model that showed sound
intermediate work but reached a wrong final answer. That bias inflates judge-audited scores for
exactly the verbose, reasoning-heavy models that strict scoring penalises. The two regimes err in
opposite directions, which is the argument for reporting both and treating the interval between them
as the honest statement of what is known.

\subsection{Scale Does Not Explain These Results}

Of the 66 pairwise comparisons, 15 survive Bonferroni correction (Section~\ref{sec:significance}).
The surviving differences separate the top of the table from the bottom; the middle is largely
undifferentiated, and two comparisons make the point sharply.

Llama 4 Scout 17B and LLaMA-3.3-70B are statistically indistinguishable (83.3\% vs 83.7\%,
$p = 0.786$) across a four-fold parameter difference. GPT-OSS 120B and GPT-OSS 20B are likewise
indistinguishable (77.1\% vs 76.8\%, $p = 0.912$) across a six-fold difference within a single model
family. The second case admits a stronger statement than equal accuracy: the two GPT-OSS models have
the highest error overlap of any pair in the study (Jaccard 0.427 against a median of 0.274 across
all 66 pairs). They do not merely score alike---they fail on substantially the same items. Whatever
limits them on this benchmark is a property the additional parameters did not address.

\subsection{What This Means for Benchmark Construction}

Two findings here were not things we set out to measure, and both concern the benchmark rather than
the models.

The first is that half our items do not discriminate. 213 of 406 are answered correctly by at least
11 of 12 models. We built the benchmark for coverage of a regulatory domain, which is a reasonable
objective, but coverage and discrimination are different targets and expert annotation effort spent
on one does not buy the other. A benchmark reporting only nominal size overstates the evidence
behind any leaderboard gap it produces.

The second concerns our difficulty labels. We assigned difficulty by counting reasoning steps, and
LLaMA-3.3-70B performs \emph{better} on our hard items than our easy ones. The likeliest reading is
that reasoning-step count captures one axis of difficulty while missing another---hard items here
tend to state their amendment-chain cues explicitly, whereas medium items turn on a single subtle
qualifier. We report the label scheme's failure rather than repairing it after the fact, because the
inversion is informative about what makes regulatory text hard.

\section{Conclusion}
\label{sec:conclusion}

Benchmark accuracy is reported as a property of a model. It is a property of a model and a scoring
rule together, and on IndiaFinBench those two contributions are separable and large. Scoring twelve
LLMs under strict string matching and under a full-coverage cross-model judge audit of every
REG/NUM/TMP prediction produces two rankings that agree only weakly: the model placed last under strict scoring places
first under judge-audited scoring, because 93.9\% of its strict errors are reclassified as correct
by the cross-model judge, primarily reflecting format differences rather than incorrect reasoning. Neither regime is wrong. They answer different questions, and reporting
one alone leaves the reader unable to tell which.

Underneath that, half the benchmark carries no discriminative signal: 213 of 406 items are answered
correctly by at least 11 of the 12 models, leaving 134 that separate systems. Nominal size and
effective size are different quantities, and the difference bounds how much weight any reported gap
can carry.

IndiaFinBench is released in full---406 expert-annotated items over 34 of 192 collected SEBI and RBI documents,
the evaluation harness, per-item judge verdicts, and all model predictions---so that both regimes
can be reproduced and both claims contested. The context-only setup isolates inference over supplied
regulatory passages from recall of memorised training text, although it cannot rule out training-data
exposure to the underlying documents; the Indian regulatory setting makes this isolation meaningful
because the source text is jurisdiction-specific and less likely to dominate pretraining corpora than
US or European financial text. Whether these scoring-regime effects generalize beyond Indian
financial regulation is an empirical question for future evaluations; our results motivate reporting
both regimes on other benchmarks.
